% =====================================================================
% Sección: Conclusiones y recomendaciones
% Responsable: Parte 7 — Daniel Alonso Gracia Pinto
% =====================================================================
\section{Conclusiones}
\label{sec:conclusiones}

Las 130 corridas permiten responder la pregunta planteada en la
sección~\ref{sec:pregunta} con evidencia estadística. Las conclusiones se resumen a
continuación.

\begin{enumerate}[label=\arabic*.]
    \item \textbf{Respuesta a la pregunta de investigación.} En la MANET jerárquica de
    dos niveles simulada, con enrutamiento AODV, la movilidad a nivel de clúster degrada
    el desempeño mucho más que la movilidad a nivel de nodo. A igual velocidad, el
    escenario con solo movilidad de clúster (E2) entrega entre 16 y 27 puntos
    porcentuales menos paquetes que el escenario con solo movilidad de nodo (E1), y su
    retardo es entre 1.3 y 4.8 veces mayor.

    \item \textbf{Se acepta la hipótesis principal.} La prueba $t$ de Welch unilateral
    rechaza $H_0$ para la PDR en las cuatro velocidades ($p \leq 0.002$) y para el
    retardo en tres de las cuatro (5, 10 y 20\,m/s). Ambas métricas cumplen la regla de
    decisión fijada antes del experimento (al menos tres de cuatro velocidades). Los
    resultados son coherentes con el mecanismo propuesto: la movilidad de clúster rompe
    los enlaces entre CH, por donde pasa todo el tráfico entre clústeres, mientras que
    la movilidad de nodo solo altera rutas locales.

    \item \textbf{El efecto de la velocidad depende del tipo de movilidad (HS1,
    parcialmente respaldada).} Con movilidad de clúster, mayor velocidad implica menor
    PDR y, sobre todo, mucho más retardo y jitter (el retardo pasa de unos 25\,ms a
    94--101\,ms entre 1 y 20\,m/s). Con solo movilidad de nodo, la velocidad no tiene
    efecto medible, porque en este diseño un miembro nunca sale del alcance de su CH.

    \item \textbf{La movilidad de clúster domina el efecto combinado (HS2).} El escenario
    con ambas movilidades (E3) no es significativamente distinto de E2 en ninguna
    velocidad: la degradación no se suma, sino que la determina el nivel superior de la
    jerarquía.

    \item \textbf{La brecha entre ambos tipos de movilidad crece con la velocidad (HS3).}
    La diferencia de retardo entre E2 y E1 crece de forma monótona (de 6\,ms a 74\,ms) y
    la de PDR crece hasta 10\,m/s y luego se estabiliza en unos 26 puntos.

    \item \textbf{Implicación práctica.} En redes jerárquicas, la estabilidad de la
    columna vertebral entre cluster heads es el factor crítico del desempeño. Los
    esfuerzos de diseño (elección de CH, alcance del backbone, rutas de respaldo entre
    clústeres) rinden más que proteger los enlaces dentro de cada clúster.
\end{enumerate}

\section{Recomendaciones}
\label{sec:recomendaciones}

\subsection{Limitaciones del estudio}

El estudio tiene las siguientes limitaciones:

\begin{itemize}
    \item \textbf{Velocidad efectiva de Random Waypoint.} Los clústeres hacen una pausa
    fija de 2\,s en cada destino, y a mayor velocidad pasan una fracción mayor del tiempo
    detenidos (cerca del 4\,\% a 1\,m/s y del 21\,\% a 20\,m/s). La velocidad media
    efectiva es, por tanto, menor que la nominal, un problema conocido de este modelo
    \cite{yoon2003harmful, navidi2004stationary}.

    \item \textbf{Geometría que protege a E1.} Con regiones de $\pm 40$\,m y alcance de
    60\,m, un miembro siempre queda al alcance directo de su CH. Esto explica la robustez
    de E1, pero también significa que la conclusión depende de la relación entre el
    tamaño de la región de los nodos y el alcance del clúster. Con regiones más grandes,
    la movilidad de nodo podría pesar más.

    \item \textbf{Modelo de propagación idealizado.} El modelo de disco no representa
    desvanecimiento, sombras ni interferencia gradual: un enlace existe por completo o no
    existe.

    \item \textbf{Escala reducida.} Se usaron 20 nodos, 4 clústeres, 8 flujos de baja
    tasa (16\,kbit/s), CH fijos dentro de su clúster y un solo protocolo (AODV). Los
    resultados no deben generalizarse a cualquier MANET sin nuevos experimentos.

    \item \textbf{Comparaciones múltiples.} Las pruebas de hipótesis se hicieron por
    velocidad con $\alpha = 0.05$ sin ajuste por comparaciones múltiples. Para la PDR
    esto no cambia la conclusión ($p \leq 0.002$ en todas las velocidades), pero la
    prueba de retardo a 5\,m/s ($p = 0.029$) no superaría una corrección de Bonferroni.

    \item \textbf{Diez réplicas por configuración.} Los intervalos de confianza de E2 y
    E3 son amplios (por ejemplo, $\pm 21$\,ms en el retardo de E2 a 20\,m/s). Más
    réplicas permitirían detectar diferencias menores, como la de retardo a 1\,m/s.
\end{itemize}

\subsection{Trabajo futuro}

Se proponen las siguientes líneas de trabajo:

\begin{itemize}
    \item Usar modelos de movilidad sin el sesgo de velocidad de Random Waypoint, como
    Gauss-Markov o Random Direction, y un modelo de grupo como RPGM
    \cite{hong1999group, camp2002survey}.
    \item Variar el tamaño de la región de los nodos respecto al alcance del clúster,
    para encontrar el punto en que la movilidad de nodo empieza a degradar la red.
    \item Comparar AODV con protocolos proactivos como OLSR \cite{clausen2003olsr} o
    DSDV, que reaccionan de forma distinta a la ruptura del backbone.
    \item Medir la sobrecarga de enrutamiento (paquetes de control de AODV por paquete
    de datos entregado), que en los datos ya disponibles crece claramente con la
    movilidad de clúster.
    \item Escalar a más nodos y niveles jerárquicos, mover también los superclústeres
    (\texttt{--superSpeed}) y permitir la reelección dinámica de CH.
    \item Explorar ns3-ai para adaptar en línea parámetros de enrutamiento o la
    elección de CH según la movilidad observada.
\end{itemize}
